\section{The adversary}\label{sec:overview}\label{sec:algorithm}

We describe the data structure maintained by our adversary and how it is
updated in response to comparison queries from the sorting algorithm.

Kaplan, Zamir, and Zwick~\cite{kzz} retain disjoint ordered pairs within a
node. Our adversary maintains a rooted tree with a DAG at each node, so
several retained comparisons can involve the same item. For example, it
can keep $a<b$ and $a<c$ together while leaving the order of $b$ and $c$
undecided. This lets it represent richer partial orders before moving
items further down the tree.

The computational challenge is to use these richer orders without having
to approximately count the linear extensions of an entire DAG. We
periodically simplify a node's DAG by removing vertices and moving their
items into child nodes, creating children when needed. Constraints at the
parent preserve the comparisons involving the moved items, while their
order within a child can remain unresolved. These simplifications allow
comparison answers to be chosen by deterministic calculations on small
neighborhoods in that node's DAG. The analysis must show that the resulting loss
of possible orders is small enough to retain the optimal comparison bound.

\subsection{The data structure}

Our data structure is a rooted tree. Each node stores item records and
a single directed acyclic graph (DAG), whose vertices are the items
recorded at that node. We call this its \emph{in-node DAG}. Figure~\ref{fig:overview-state} shows an example.

Each item has a record $(x,X_x,\ell_x,r_x)$ stored at exactly one node.
The label $x$ identifies its key, and $X_x$ is a symbolic variable for
its position among the items at that node. The fields $\ell_x,r_x$ restrict
this variable to the open interval $(\ell_x,r_x)\subseteq(0,1)$.
The bounds $\ell_x,r_x$ are fractions whose denominators are powers of two;
we call such numbers \emph{dyadic rationals}.
We call an item \emph{free} at its node while its position remains variable.

In an in-node DAG, an arc $x\to y$ records a comparison answered $x<y$
between two items stored there, imposing $X_x<X_y$ on their positions.
Each answered comparison between free items contributes at most one arc.

A \emph{feasible assignment} at a node is a choice of real values for its
position variables satisfying $\ell_x<X_x<r_x$ for every item $x$ and
$X_x<X_y$ for every arc $x\to y$ in its in-node DAG.

We keep each stored interval equal to the range of values its coordinate
can take in feasible assignments. An arc $x\to y$ propagates lower bounds
from $x$ to $y$ and upper bounds from $y$ to $x$. After changing a node,
the routine \texttt{Normalize} makes one pass in topological order for the
lower bounds and one in reverse order for the upper bounds.
Algorithm~\ref{alg:normalize} gives the two passes; Appendix~\ref{app:geometry}
proves that they preserve exactly the same feasible assignments.

Initially the tree consists of just the root, containing all records
with bounds $\ell_x=0$ and $r_x=1$. The position variables $X_x$ are
unassigned, and the root's in-node DAG has no arcs.
Thus every assignment in $(0,1)^n$ is feasible, and every total order
can be obtained by assigning distinct positions.

Each child is attached to a fixed position $c\in(0,1)$ in its parent,
called its \emph{anchor}. A node has at most one child at each anchor.
For any feasible assignment at the parent, an
item with position below $c$ precedes every item in that child's subtree;
an item with position above $c$ follows every item in it. Subtrees of
different children are ordered by their anchors. All items of each child
subtree therefore occur consecutively in the resulting total order.
Within the child, position variables again range over $(0,1)$, and its
in-node DAG and interval bounds describe the possible orders inside this block.

\begin{center}
\begin{minipage}{\linewidth}
\captionsetup{hypcap=false}
\centering
\begin{tikzpicture}[>=Latex,font=\small,
  group/.style={draw=ink!65,rounded corners=2pt,fill=pale},
  item/.style={circle,draw=ink,fill=white,minimum size=6mm,inner sep=1pt}]
% Root: the graph, the interval bounds, and the child anchors.
\draw[group] (-6.8,0) rectangle (6.8,4.7);
\node[font=\bfseries] at (0,4.35) {Root};
\node at (-3.8,3.85) {In-node DAG};
\node at (3.2,3.85) {Allowed intervals};
\draw[ink!20] (-.5,1.85)--(-.5,4.05);
\node[item] (oa) at (-5,2.7) {$a$};
\node[item] (ob) at (-2.5,3.3) {$b$};
\node[item] (oc) at (-2.5,2.1) {$c$};
\draw[->,thick,accent] (oa)--(ob);
\draw[->,thick,accent] (oa)--(oc);
\foreach \x/\label in {1/0,2.15/{1/4},3.3/{1/2},4.45/{3/4},5.6/1}
  \node[font=\scriptsize] at (\x,3.45) {$\label$};
\foreach \y/\label/\lo/\hi in {3.05/a/1/4.45,2.55/b/2.15/5.6,2.05/c/2.15/5.6}{
  \node[left] at (.65,\y) {$X_{\label}$};
  \draw[ink!25] (1,\y)--(5.6,\y);
  \draw[accent,very thick] (\lo,\y)--(\hi,\y);
  \draw[accent,fill=white] (\lo,\y) circle (1.7pt);
  \draw[accent,fill=white] (\hi,\y) circle (1.7pt);
}
\node at (0,1.6) {Child anchors in the root's unit interval};
\draw[thick] (-5.8,.7)--(5.8,.7);
\node[below] at (-5.8,.65) {$0$};
\node[below] at (5.8,.65) {$1$};
\fill (-2.9,.7) circle (2.3pt);
\fill (2.9,.7) circle (2.3pt);
\node[above=3pt] at (-2.9,.7) {$1/4$};
\node[above=3pt] at (2.9,.7) {$3/4$};
\draw[thick,ink!55] (-2.9,.7)--(-2.9,-.2)--(-3.55,-.9);
\draw[thick,ink!55] (2.9,.7)--(2.9,-.2)--(3.55,-.9);
% Each child repeats both kinds of information on its own unit interval.
\begin{scope}[shift={(-3.55,-.9)}]
\draw[group] (-3.25,-3) rectangle (3.25,0);
\node[font=\bfseries] at (0,-.35) {Child at $1/4$};
\node[font=\footnotesize] at (-1.9,-.85) {In-node DAG};
\node[font=\footnotesize] at (1.5,-.85) {Allowed intervals};
\draw[ink!20] (-.55,-1.1)--(-.55,-2.5);
\node[item] (op) at (-2.55,-1.7) {$p$};
\node[item] (oq) at (-1.25,-1.7) {$q$};
\draw[->,thick,accent] (op)--(oq);
\foreach \x/\label in {.4/0,2.8/1}
  \node[font=\scriptsize] at (\x,-1.25) {$\label$};
\foreach \y/\label/\lo/\hi in {-1.7/p/.4/2.8,-2.2/q/.4/2.8}{
  \node[left] at (.15,\y) {$X_{\label}$};
  \draw[ink!25] (.4,\y)--(2.8,\y);
  \draw[accent,very thick] (\lo,\y)--(\hi,\y);
  \draw[accent,fill=white] (\lo,\y) circle (1.7pt);
  \draw[accent,fill=white] (\hi,\y) circle (1.7pt);
}
\node[font=\footnotesize] at (0,-2.7) {Positions within this child};
\end{scope}
\begin{scope}[shift={(3.55,-.9)}]
\draw[group] (-3.25,-3) rectangle (3.25,0);
\node[font=\bfseries] at (0,-.35) {Child at $3/4$};
\node[font=\footnotesize] at (-1.9,-.85) {In-node DAG};
\node[font=\footnotesize] at (1.5,-.85) {Allowed intervals};
\draw[ink!20] (-.55,-1.1)--(-.55,-2.5);
\node[item] (orr) at (-2.55,-1.7) {$r$};
\node[item] (os) at (-1.25,-1.7) {$s$};
\draw[->,thick,accent] (orr)--(os);
\foreach \x/\label in {.4/0,2.8/1}
  \node[font=\scriptsize] at (\x,-1.25) {$\label$};
\foreach \y/\label in {-1.7/r,-2.2/s}{
  \node[left] at (.15,\y) {$X_{\label}$};
  \draw[accent,very thick] (.4,\y)--(2.8,\y);
  \draw[accent,fill=white] (.4,\y) circle (1.7pt);
  \draw[accent,fill=white] (2.8,\y) circle (1.7pt);
}
\node[font=\footnotesize] at (0,-2.7) {Positions within this child};
\end{scope}
\end{tikzpicture}
\captionof{figure}{A state of the adversary. Each large box is a tree node,
containing an in-node DAG and one open interval for each of its items.}
\label{fig:overview-state}
\end{minipage}
\end{center}

In the root of Figure~\ref{fig:overview-state}, the in-node DAG records $a<b$ and
$a<c$.
To read an order from this state, choose feasible positions at each node,
avoiding ties with other items or child anchors. List the items and anchors
from left to right, and recursively replace each anchor by its child's
order. For example, root positions
$(X_a,X_b,X_c)=(1/3,1/2,7/8)$ give
\[
p<q<a<b<r<s<c.
\]
Other feasible choices give other orders. Every order represented this way
must respect all previous answers to the sorter. We call the sequence of
these answers the \emph{transcript}.

\subsection{Answering a comparison}\label{sec:comparison}

Algorithm~\ref{alg:compare} shows the algorithm the adversary uses to answer
a comparison by the sorter.

To compare $A$ and $B$, start at the root and descend into a child whenever
its subtree contains both items. Stop when at least one item is stored at
the current node, or when the items lie in different child subtrees.
At this node, use an item's free coordinate if it is stored here, and the
corresponding child anchor otherwise.
In Figure~\ref{fig:overview-state}, comparing $p$ with $q$ descends into
the left child; comparing $p$ with $r$ uses the anchors $1/4$ and $3/4$;
and comparing $a$ with $p$ compares $X_a$ with $1/4$ at the root.

Two distinct anchors determine the answer immediately.

For a query involving a free coordinate, consider the uniform distribution on feasible
positions at this node. The subroutine
\texttt{Certify} computes two nonnegative rational numbers $p_{AB},p_{BA}$
that are lower bounds on the probabilities of $A<B$ and $B<A$.
Fix a rational tolerance $0<\eta<1/32$. If the larger certified fraction
is at least $1/2-\eta$, choose that answer, breaking a tie in favor of
$A<B$. Otherwise, move a queried free item into a child and restart the
query. We call this move a \emph{refinement}, denoted \texttt{Refine}$(x,\mathit{reason})$;
Algorithm~\ref{alg:refine} in Section~\ref{sec:overview-refine} gives the
operation. Section~\ref{sec:potential} shows why a failed certificate
makes such a refinement affordable.
For two free items, the events are $X_A<X_B$ and $X_B<X_A$; for a free
item and an anchor, they compare its coordinate with the fixed anchor.
An answer between two free items adds an arc to the in-node DAG;
an answer involving an anchor adds an interval bound.

After recording an answer, we repair its node if the new arc has made
the in-node DAG too large. Fix integers $k,H\ge2$. Between completed
calls, every vertex has fewer than $k$ neighbors and every directed
path has fewer than $H$ vertices. Degree counts both incoming and
outgoing neighbors, and path length counts vertices.
If a path of $H$ vertices has appeared, refine all its items in path
order. Otherwise, if a vertex has reached degree $k$, refine that vertex.
Algorithm~\ref{alg:repair} gives this operation. One such repair restores
both bounds: any new violation must involve the arc just inserted, and
the repair removes that arc. Section~\ref{sec:repairs} proves this claim.

The following lemma supplies the guarantee needed by the comparison rule.
Section~\ref{sec:oracle} implements \texttt{Certify} by counting finite
calculations in which enclosing intervals establish one of the two answers.

\begin{lemma}[Certified comparison probabilities]\label{lem:estimation}
Suppose a node admits a feasible assignment, and its in-node DAG has
degree at most $k$ and directed paths of at most $2H$ vertices.
For two distinct free coordinates, or a free coordinate and
a dyadic threshold, \textnormal{\texttt{Certify}} deterministically returns
nonnegative rational lower bounds on the probabilities of the two answers.
\end{lemma}

The recursive construction is needed for at most
\begin{equation}\label{eq:comparison-horizon}
Q(n)=n\ceil{\log n}
\end{equation}
answers. Record these answers in a list $\mathcal T$, initially empty.
If another query arrives after $Q(n)$ answers, choose a fixed total order
consistent with $\mathcal T$ and use it for every subsequent answer.
Section~\ref{sec:runtime} constructs this order by a deterministic
topological sort. A sorter reaching this point has already made at least
$n\log n$ comparisons.

\begin{algorithm}[H]
\caption{\texttt{Compare}$(A,B)$: answer a comparison query.}
\label{alg:overview-compare}\label{alg:compare}
\DontPrintSemicolon
\KwIn{Distinct item labels $A,B$ and the current adversary state.}
\KwOut{An answer $A<B$ or $B<A$, with the state updated.}
\If{fixed ranks have been chosen}{\Return their order of $A,B$\;}
\If{$|\mathcal T|=Q(n)$}{
  Form the DAG of answers in $\mathcal T$ and choose a topological order, breaking ties by label\;
  Store its ranks and \Return their order of $A,B$\;
}
\While{true}{
  $v\gets$ the root\;
  \While{$A$ and $B$ lie in the same child subtree of $v$}{
    $v\gets$ that child\;
  }
  \If{$A$ and $B$ lie in different child subtrees of $v$}{
    $\mathit{answer}\gets$ the left-to-right order of their child anchors\;
    \textbf{break}\;
  }
  $(p_{AB},p_{BA})\gets$ \texttt{Certify}$(v,A,B,\eta)$\;
  \If{$\max(p_{AB},p_{BA})\ge1/2-\eta$}{
    $\mathit{answer}\gets A<B$ if $p_{AB}\ge p_{BA}$, and $B<A$ otherwise\;
    \textbf{break}\;
  }
  $x\gets$ the selected queried free item\;
  \texttt{Refine}$(x,\mathrm{failure})$\;
}
\If{at least one queried item is stored at $v$}{
  Store $\mathit{answer}$ as an arc in the in-node DAG at $v$ or as an interval bound\;
  \texttt{Normalize}$(v)$\;
}
Append $\mathit{answer}$ to $\mathcal T$\;
\texttt{Repair}$(v)$\;
\Return $\mathit{answer}$\;
\end{algorithm}

Before recording the answer, only the two queried items can move.
Section~\ref{sec:repairs} shows that each moves at most once in this
part of the call. Thus the loop makes at most two refinements and three
comparison-certificate calculations. A queried item may move again in
the subsequent repair; the whole call makes at most $n+2$ refinements.

For deterministic selection, choose the smallest eligible item label and
the lexicographically first eligible directed path. Complete the entire
selected path batch, even if earlier refinements have already removed
its arcs.

\begin{algorithm}[H]
\caption{\texttt{Repair}$(v)$: restore the in-node DAG bounds after an answer.}
\label{alg:repair}
\DontPrintSemicolon
\KwIn{The node $v$ at which the answer was determined.}
\If{the in-node DAG at $v$ has a directed path on $H$ vertices}{
  $L\gets$ the selected path's list of labels, in path order\;
  \ForEach{$x\in L$ in path order}{\texttt{Refine}$(x,\mathrm{repair})$\;}
}
\ElseIf{a vertex at $v$ has degree $k$}{
  $x\gets$ the selected such vertex\;
  \texttt{Refine}$(x,\mathrm{repair})$\;
}
\end{algorithm}

\subsection{Moving an item into a child}\label{sec:overview-refine}

Algorithm~\ref{alg:refine} gives \texttt{Refine}$(x,\mathit{reason})$,
which moves $x$ from its current node into a child. The argument
$\mathit{reason}$ records whether the move follows a failed comparison
certificate or is part of a structural repair.

After a failed comparison certificate, choose an anchor in the middle
half of $(\ell_x,r_x)$. For a structural repair, first try to find an
interval with substantial probability inside and on both sides. If that
calculation fails, use the middle half instead. In either case, a simple
rounding rule selects a dyadic position in the chosen interval;
Algorithm~\ref{alg:choose-anchor} in Section~\ref{sec:median-search}
gives the rule.

\Needspace{6\baselineskip}
\begin{lemma}[A central interval]\label{lem:central-interval}
Under the hypotheses of Lemma~\ref{lem:estimation},
\textnormal{\texttt{CentralInterval}} deterministically returns either
failure or dyadic $u<v$ in $(\ell_x,r_x)$ satisfying
\[
\Prb(X_x\le u)\ge\frac14,\qquad
\Prb(u\le X_x\le v)\ge\frac14,\qquad
\Prb(X_x\ge v)\ge\frac14.
\]
\end{lemma}

Section~\ref{sec:oracle} implements this routine using the same enclosing
intervals as \texttt{Certify}.

Fixing $X_x=c$ places $x$ in the child block anchored at $c$.
The arcs incident to $x$ in that in-node DAG become interval bounds: $y\to x$ gives $X_y<c$,
and $x\to z$ gives $X_z>c$. Remove $x$ and these arcs from the parent,
and retain the new bounds on its neighbors.
Insert $x$ into the child anchored at $c$, creating it if needed. In that
child, $x$ receives a fresh position variable with interval $(0,1)$ and no incident arcs; the
child's existing constraints remain unchanged.
Figure~\ref{fig:overview-refine} illustrates reuse of an existing child.

\begin{algorithm}[H]
\caption{\texttt{Refine}$(x,\mathit{reason})$: move an item into a child.}
\label{alg:refine}
\DontPrintSemicolon
\KwIn{A free item $x$, the reason for its move, and the current adversary state.}
\KwOut{The updated state, with $x$ stored in a child of its former node.}
$v\gets$ the node storing $x$\;
$c\gets$ \texttt{ChooseAnchor}$(x,\mathit{reason})$\;
\ForEach{arc $y\to x$ in the in-node DAG at $v$}{
  $r_y\gets\min(r_y,c)$\;
}
\ForEach{arc $x\to z$ in the in-node DAG at $v$}{
  $\ell_z\gets\max(\ell_z,c)$\;
}
Remove $x$'s record from $v$ and its vertex and incident arcs from the in-node DAG at $v$\;
\texttt{Normalize}$(v)$\;
\If{$v$ has no child anchored at $c$}{
  Create a child of $v$ at $c$, with no records and an empty in-node DAG\;
}
$u\gets$ the child of $v$ anchored at $c$\;
Insert a record $(x,X_x,0,1)$ at $u$, with a fresh, unassigned position variable $X_x$\;
Insert $x$ as an isolated vertex in the in-node DAG at $u$\;
\end{algorithm}

The propagation routine used above keeps the stored intervals up to date.
For example, tightening $r_y$ also tightens the upper bound of every
predecessor of $y$. The two passes below perform all such propagations.

\begin{algorithm}[H]
\caption{\texttt{Normalize}$(v)$: propagate interval bounds through the in-node DAG.}
\label{alg:normalize}
\DontPrintSemicolon
\KwIn{A node $v$ whose arcs or interval bounds have changed.}
\ForEach{vertex $x$ in topological order}{
  $\ell_x\gets\max(\{\ell_x\}\cup\{\ell_y:y\to x\})$\;
}
\ForEach{vertex $x$ in reverse topological order}{
  $r_x\gets\min(\{r_x\}\cup\{r_y:x\to y\})$\;
}
\end{algorithm}

\begin{center}
\begin{minipage}{\linewidth}
\captionsetup{hypcap=false}
\centering
\begin{tikzpicture}[>=Latex,font=\small,
  group/.style={draw=ink!65,rounded corners=2pt,fill=pale},
  item/.style={circle,draw=ink,fill=white,minimum size=6mm,inner sep=1pt}]
\begin{scope}[shift={(-3.2,0)}]
\node at (0,1.55) {Before refinement};
\draw[group] (-2.4,-.45) rectangle (2.4,1.2);
\node[item] (ra) at (-1.5,.65) {$a$};
\node[item] (rx) at (0,.65) {$x$};
\node[item] (rb) at (1.5,.65) {$b$};
\draw[->,thick,accent] (ra)--(rx);
\draw[->,thick,accent] (rx)--(rb);
\fill (0,-.15) circle (2pt) node[right=4pt] {anchor $c$};
\draw[thick,ink!55] (0,-.15)--(0,-1);
\draw[group] (-1.8,-2) rectangle (1.8,-1);
\node at (0,-1.3) {Existing child at $c$};
\node at (0,-1.7) {$z$ free};
\end{scope}
\draw[->,thick] (-.55,-.2)--(.55,-.2);
\begin{scope}[shift={(3.2,0)}]
\node at (0,1.55) {After fixing $X_x=c$};
\draw[group] (-2.4,-.45) rectangle (2.4,1.2);
\node[item] at (-1.5,.65) {$a$};
\node[item] at (1.5,.65) {$b$};
\node at (0,.65) {$X_a<c<X_b$};
\fill (0,-.15) circle (2pt) node[right=4pt] {anchor $c$};
\draw[thick,ink!55] (0,-.15)--(0,-1);
\draw[group] (-1.8,-2) rectangle (1.8,-1);
\node at (0,-1.3) {Same child at $c$};
\node at (0,-1.7) {$x,z$ free; no arc};
\end{scope}
\end{tikzpicture}
\captionof{figure}{The two arc inequalities become bounds on the remaining parent
coordinates. The new coordinate of $x$ orders it within the child, where
both $x<z$ and $z<x$ remain possible. The anchor $c$ locates the child
block among the parent items; the child's constraints determine the orders
within that block.}
\label{fig:overview-refine}
\end{minipage}
\end{center}

Thus the parent fixes where the child block occurs, while the order of
$x$ within that block stays open. Lemma~\ref{lem:admission} proves that
this operation preserves all answered comparisons.

Refinement can impose additional order relations. For example, suppose
the root contains $A,B,C$ with only $A<B$ recorded. Refining $A$ at $1/2$
restricts $B$ to $(1/2,1)$. Refining $C$ into the same child then leaves
both $A<C<B$ and $C<A<B$ possible, while fixing the additional relation
$C<B$. Section~\ref{sec:potential}
shows that we can afford this loss of possible orders while retaining the
optimal asymptotic comparison bound.

Table~\ref{tab:parameters} collects the thresholds used by
Algorithms~\ref{alg:compare}--\ref{alg:normalize} above.
In the next section, we explain why the adversary forces any correct
sorter to perform many comparisons.
